% ---
% title: The Gradiant of Reward Function
% date: 2026-02-17
% tag: RL
% excerpt: Why is it so important to get the gradient of the reward function?
% ---
\section{问题}

为什么得到奖励函数的梯度这么重要呢？

这句话在说一种“更强的假设”：\textbf{奖励 \(r(x)\) 对最终生成图像 \(x\) 的梯度是可用的}\(\nabla_x r(x)\) 。一旦有了它，你就能把奖励当成一个可微的“损失头”，把梯度从最终图像\textbf{直接反传到扩散模型参数}，训练会变得又快又稳，所以很多工作喜欢这么设定。

为什么它这么重要？核心在这几件事：

\subsection{有奖励梯度 = 低方差、可直接优化}

如果 \(\nabla_x r(x) \)可得，你可以写成

\[
\max_\theta \;\mathbb{E}_{x \sim p_\theta}[r(x)] \quad \Rightarrow \quad \min_\theta \; \mathbb{E}[ -r(x)]
\]

因为 x 是由扩散采样链 \(x = g_\theta(\epsilon)\) 生成的（可微计算图），所以

\[
\nabla_\theta \mathbb{E} [r(g_\theta(\epsilon))] = \mathbb{E} \left[ \underbrace{\nabla_x r(x)}_{\text{reward梯度}} \cdot \underbrace{\nabla_\theta g_\theta(\epsilon)}_{\text{生成过程对参数的梯度}} \right]
\]

这就是\textbf{路径导数（pathwise gradient）/重参数化梯度}：不需要对随机采样分布做 REINFORCE 式的 \(\nabla_\theta \log p_\theta(x) \)估计，通常\textbf{方差更小、收敛更快}。

\hrule

\subsection{没有奖励梯度就得走“RL/score-function”路线，方差大、难稳定}

如果拿不到 \(\nabla_x r(x)\)，你仍想优化\( \mathbb{E}[r(x)]\)，典型就只能用：

\[
\nabla_\theta \mathbb{E}_{x \sim p_\theta}[r(x)] = \mathbb{E}[r(x)\nabla_\theta \log p_\theta(x)]
\]

这类估计（score-function / REINFORCE）在扩散里非常痛苦，因为：

\begin{itemize}
  \item x 是一条\textbf{很长的随机链}的结果（几十步、上百步），\(\log p_\theta(x)\) 的梯度噪声大
  \item reward 通常只在最终图像给一次信号，\textbf{信用分配（credit assignment）}极难
  \item 需要 baseline、value function、advantage 等来降方差（你之前问的 value function 就是为这个服务）
\end{itemize}

所以论文里常说 “high-variance Monte Carlo estimators”。

\hrule

\subsection{有奖励梯度可以做“终态反传”，实现极简、样本效率高}

你看到的那句 “directly backpropagate gradient signals from the final state” 指的是：

\begin{enumerate}
  \item 采样得到最终图像 \(x_0\)
  \item 计算 reward（可微）\(r(x_0)\)
  \item 把 \(-r(x_0)\) 当 \(loss\)，沿着采样计算图（或某种可微近似）\textbf{直接反传到 UNet/LoRA}
\end{enumerate}

这样做的工程优势特别大：不需要 reward model 的训练、不需要 value function、不需要复杂的 RL 算法结构。

\hrule

\subsection{但“奖励梯度可得”在现实里往往不成立}

这也是为什么很多方法强调它是“假设”：

\begin{itemize}
  \item 人类偏好、人工打分、A/B 选择：\textbf{不可微}
  \item 外部评价器（比如黑盒 API、审美模型只返回分数但不给梯度）：\textbf{黑盒}
  \item 许多评价指标（FID、CLIP 变体里带非光滑操作、离散规则/安全策略）：可能\textbf{不可微或梯度不可靠}
  \item 即便可微，也可能存在\textbf{梯度 hacking / reward hacking}：模型学会骗 reward 的梯度而非提升真实质量
\end{itemize}

所以一些工作转向：只用 reward 标量 + KL 约束（你 SQDF 的 KL 项就是为了稳定），或者用 DPO/GRPO/偏好学习等。

\hrule

\subsection{和你 SQDF 这类方法的关系}

你现在的训练里，reward（审美）是通过 online\_loss\_fn(img) 得到的标量，然后你用

\[
\mathcal{L} = -\left(\gamma^{t}\, r(x) - \alpha\, KL\right)
\]

你并\textbf{不依赖} \(\nabla_x r(x)\)（至少从训练逻辑上你把 reward 当成标量权重），而是用 KL 项把更新限制在“像 ref 的一步转移附近”，用 replay/buffer 提高数据利用率。

如果 reward 梯度真的可得，你还可以进一步把 reward 对 image 的梯度直接回传（更像“纯 differentiable objective”），但你就必须确保 reward 模块是可微且梯度可信。

\hrule

一句话总结：\textbf{奖励函数梯度重要，是因为它把“生成模型对人类偏好的优化”从高方差 RL 变成低方差、可直接反传的端到端学习；但这个假设在真实 reward 场景里通常不成立，所以才会出现 SQDF/GRPO/DDPO 这种不依赖 reward 梯度的路线。}

\section{\(\nabla_x r(x)\) 怎么求}

\(\underbrace{\nabla_x r(x)}_{\text{reward梯度}}\) 的“求法”取决于你的 \textbf{reward} r(x) 是不是对输入 \(x \)可微。

下面按最常见的 3 种情况说清楚，并给你 PyTorch 里怎么拿梯度。

\hrule

\subsection{1)\textbf{\(r(x)\)}是可微模型（最常见：CLIP/审美打分网络/RewardNet）}

假设 reward 是一个神经网络 \(r_\phi(x)\)，输入是图像张量 \(x\in\mathbb{R}^{B\times C\times H\times W}\)，输出标量（或每张图一个分数）。那 \(\nabla_x r(x)\) 就是 \textbf{对输入张量的反向梯度}，用 \(autograd\) 直接拿：

\begin{lstlisting}[language=python]
# x: (B,C,H,W) float tensor, requires_grad=True
x = x.clone().detach().requires_grad_(True)

score = reward_model(x)          # (B,) or (B,1) or scalar
score_sum = score.sum()          # 变成标量方便求导

grad_x = torch.autograd.grad(
    outputs=score_sum,
    inputs=x,
    create_graph=False,          # 只需要一阶梯度通常 False
    retain_graph=False,
    only_inputs=True
)[0]                             # grad_x shape == x shape
\end{lstlisting}

\textbf{关键点（很多人卡在这里）：}

\begin{itemize}
  \item \(x\) 必须 requires\_grad\_(True)。
  \item 计算 score 时不要包 torch.no\_grad()，不要对 \(x\) 做 .detach()。
  \item reward 模型内部如果有不可微操作（硬阈值、argmax、离散采样），梯度要么为 0 要么不稳定。
\end{itemize}

\hrule

\subsection{2)\textbf{\(r(x)\)}是一个明确公式（可解析）}

比如你定义 reward 为某种可微的函数：

\begin{itemize}
  \item \(r(x)= -\|x-x^*\|_2^2\)
  \item \(r(x)=\text{mean}(\text{softplus}(x))\)
  \item 或一些可微的频域/梯度能量
\end{itemize}

那 \(\nabla_x r(x)\) 可以手推，也可以同样交给 autograd（推荐）：

\begin{lstlisting}[language=python]
x = x.requires_grad_(True)
r = -(x - x_target).pow(2).mean()
grad_x = torch.autograd.grad(r, x)[0]
\end{lstlisting}

\hrule

\subsection{3)\textbf{\(r(x)\)}是黑盒/不可微（人类打分、API、规则指标、FID…）}

这时严格意义上 \(\nabla_x r(x)\) \textbf{不存在或拿不到}。常见替代路线：

\subsubsection{(a) 有限差分近似（能用但非常贵）}

\[
\frac{\partial r}{\partial x_i}\approx \frac{r(x+\epsilon e_i)-r(x-\epsilon e_i)}{2\epsilon}
\]

维度太高（图像像素百万级）几乎不可行，除非做随机方向估计（SPSA/ES），但噪声大、样本开销巨大。

\subsubsection{\textbf{(b) 训练一个可微 surrogate reward} \(\hat r_\phi(x)\)}

用黑盒打分数据 \((x, r(x))\) 拟合一个网络 \(\hat r\)，然后用它的梯度当 \(\nabla_x r(x)\) 近似。

\subsubsection{(c) 不用\textbf{\(\nabla_x r(x)\)}，改用 score-function/RL（REINFORCE/DDPO/GRPO/SQDF 类）}

也就是你现在这条路线：reward 只当标量信号，不要求对图像可微。

\hrule

\subsection{你在 diffusion 里“用上” \(\nabla_x r(x)\) \textbf{还需要注意的一点}}

你问的是 \(\nabla_x r(x)\)（对最终图像的梯度）。但如果你想更新 UNet 参数 \(\theta\)，还需要链式法则：

\[
\nabla_\theta r(x_0)=\nabla_{x_0} r(x_0)\cdot \frac{\partial x_0}{\partial \theta}
\]

这要求你\textbf{保留从} \(\theta\) \textbf{到} \(x_0\) \textbf{的计算图}（采样过程别 no\_grad，DDIM/ODE step 要可微）。这在扩散里会非常吃显存，所以通常要：

\begin{itemize}
  \item 少步数 / 截断反传（只反传最后几步）
  \item gradient checkpoint
  \item 或像你 SQDF 那样只做“一步”更新、避免全链反传
\end{itemize}

\hrule

如果你把你现在 online\_loss\_fn 的实现（或你用的 aesthetic reward 模型那块）贴出来，我可以按你的实际张量形状/归一化范围（[-1,1] 还是 [0,1]）给你写一段\textbf{直接返回 reward 和 \(grad_x\)}的版本，保证能接到你现在的训练里。

\section{最大化奖励\(\nabla_\theta r(x_0)\)}

你这句链式法则的直觉没问题：\textbf{如果你能把采样过程写成可微映射}

\[
x_0 = g_\theta(\epsilon,\text{prompt})
\]

那最大化期望奖励的“精确方向”（更准确说是对 \textbf{\(\mathbb{E}_\epsilon[r(g_\theta(\epsilon))]\)}的无偏梯度估计）确实可以用\textbf{路径导数}：

\[
\nabla_\theta r(x_0)=\nabla_{x_0} r(x_0)\,\cdot \frac{\partial x_0}{\partial \theta}
\]

==而且在实现里你并不会显式算巨大 Jacobian \(\frac{\partial x_0}{\partial\theta}\)，autograd 实际算的是 \(VJP := (\nabla_{x_0} r)\cdot J_{x_0}(\theta)\)。==

但困难恰恰在于：\(\frac{\partial x_0}{\partial \theta}\) \textbf{“数学上存在/框架能给”} ≠ \textbf{“工程上容易/梯度好用”}。主要卡点一般有这些：

\hrule

\subsection{采样是“长链展开”，反传代价爆炸（时间+显存）}

扩散采样是 T 次迭代（常见 20/50/100 步），每一步都跑一次（甚至两次）UNet：

\begin{itemize}
  \item \textbf{CFG} 下每步要算 cond/uncond 两次 UNet → 计算直接翻倍
  \item 反向传播需要保存每步的中间激活（或 checkpoint 重新算）
  \item 结果是：\textbf{显存非常吃紧、速度非常慢}（比普通 LoRA 训练还重很多）
\end{itemize}

所以很多工作才会做：截断反传（只反最后几步）、单步近似（你 SQDF 这种）、或者用 consistency model 缩短链。

\hrule

\subsection{采样里有随机性：你优化的是期望，但单条轨迹梯度噪声很大}

采样里常见形式：

\[
x_{t-1} = \mu_\theta(x_t,t) + \sigma_t\,\eta,\quad \eta\sim \mathcal{N}(0,I)
\]

这仍然可以重参数化（\(\eta\) 独立于 \(\theta\)），所以\textbf{理论上可微}；但问题是：

\begin{itemize}
  \item 你真正想优化的是 \(\nabla_\theta \mathbb{E}_\eta[r(x_0)]\)
  \item 单次采样得到的梯度是一个 Monte Carlo 估计，\textbf{方差很大}
  \item 想稳定就要多采样/用控制变量/用 replay buffer——又进一步加大算力需求
\end{itemize}

（你现在的 buffer、PER，其实就是在缓解这种“信号稀疏 + 方差大”。）

\hrule

\subsection{“可微 ≠ 梯度好用”：容易出现梯度消失/爆炸/不稳定}

采样链里有大量会影响梯度质量的东西，例如：

\begin{itemize}
  \item scheduler 的一些操作（clamp/thresholding 等）可能让梯度在饱和区接近 0
  \item VAE decode、图像归一化、reward 模型的归一化等，都会让尺度非常敏感
  \item 长链乘起来的 Jacobian，容易造成 \textbf{梯度爆炸或消失}（类似 RNN 长序列反传问题）
\end{itemize}

结果就是：就算 autograd 能给 \(\partial x_0/\partial\theta\)，也可能\textbf{数值上很难训练}。

\hrule

\subsection{reward 梯度本身可能“不可靠”或“被钻空子”}

即便你能拿到\( \nabla_{x_0} r(x_0)\)，也还会遇到：

\begin{itemize}
  \item reward model 的梯度可能指向“投机取巧”的方向（reward hacking）
  \item reward 模型在分布外的梯度可能乱（尤其模型被你更新后，分布迅速漂移）
  \item 人类偏好/安全规则等 reward 往往是黑盒或非光滑的，梯度本来就不存在
\end{itemize}

所以很多方法宁愿用 KL 约束、偏好对比、分布约束等来稳住。

\hrule

\subsection{你以为的“\textbf{\(\partial x_0/\partial \theta\)}很容易”常常被代码结构挡住}

像你贴的代码里，rollout 经常包着 torch.no\_grad()、或者把中间 latents .detach() 存 buffer，这些都会让：

\begin{itemize}
  \item 你拿到的 \(x_0\) \textbf{不在计算图上}
  \item autograd 根本没法从 reward 反传回 UNet 参数
\end{itemize}

你要做“终态反传”，必须把采样过程改成：\textbf{至少最后若干步在 enable\_grad 下、且不 detach}。这又回到第 1 点的显存/速度问题。

\hrule

\subsubsection{小结}

\(\frac{\partial x_0}{\partial\theta} \)在 PyTorch 里“能算”，但难点不在“能不能求导”，而在：

\begin{enumerate}
  \item \textbf{长链反传成本极高}
  \item \textbf{单样本梯度方差大}（你要的是期望最优，不是某个噪声 realization）
  \item \textbf{梯度数值不稳定}（clamp/尺度/链式乘积）
  \item \textbf{reward 梯度可能不可信}（hacking/分布外）
  \item \textbf{实际代码经常为了省显存把图切断}
\end{enumerate}

如果你愿意，我可以结合你当前 SQDF 代码结构，给你一个“折中方案”的落地改法：比如\textbf{只对最后 1～k 步开梯度}、或把 consistency model 当作可微“短路解码器”，让你能真正用到 \(\nabla_{x_0} r(x_0)\) 同时不至于显存爆掉。你更想做哪一种：最后 1 步、最后 5 步、还是 consistency 单步？

\section{为什么不能直接使用雅可比？}

你抓到关键点了：\textbf{我们从来不显式构造那个巨大 Jacobian}

\[
J_{x_0}(\theta)=\frac{\partial x_0}{\partial \theta}
\]

因为它的尺寸离谱：如果 \(x_0\) 有 \(D\) 个像素维度、\(\theta\) 有 \(P\) 个参数，那 Jacobian 是 \(D\times P\)。\(D\sim 10^5\!-\!10^6\)，\(P\sim 10^8\)，根本存不下也算不动。

但你真正需要的其实是最终梯度\(\nabla_\theta r(x_0)\)

链式法则写成矩阵形式就是：

\[
\nabla_\theta r(x_0)=\underbrace{\left(\nabla_{x_0} r(x_0)\right)}_{1\times D}\;\underbrace{\frac{\partial x_0}{\partial \theta}}_{D\times P}
\]

这是一个\( 1\times P \)的向量（对每个参数的梯度）。

这里出现的乘法

\[
(\nabla_{x_0} r)\cdot J_{x_0}(\theta)
\]

就是 \textbf{VJP（Vector–Jacobian Product，向量-雅可比乘积）}：

\begin{itemize}
  \item “vector” 指的是 \(v=\nabla_{x_0} r(x_0)\)（形状 \(D\)）
  \item “Jacobian” 指的是 \(J=\partial x_0/\partial\theta\)（形状 \(D\times P\)）
  \item “product” 得到 \(v^\top J\)（形状 P），也就是你要的 \(\nabla_\theta r\)
\end{itemize}

\subsection{为什么 autograd 说它在算 VJP？}

反向传播（reverse-mode autodiff）的核心能力就是：\textbf{给定输出的梯度（上游梯度），把它往前“乘”每一层的 Jacobian 的转置}。整体效果等价于一次 VJP。

在 PyTorch 里，假设：

\begin{lstlisting}[language=python]
x0 = g_theta(...)          # x0: shape (D,)
r  = reward(x0)            # r: scalar
r.backward()
\end{lstlisting}

backward() 做的就是把标量梯度 \(dr/dr=1\) 作为起点，一路往回传。到达 \(x_0\) 这一层时，上游梯度就是 \(\nabla_{x_0} r\)，然后继续往回把它乘上每一步的局部 Jacobian，最终得到 \(\nabla_\theta r\)。\textbf{你从未显式生成} \(J\)，但得到了 \(v^\top J\)。

\hrule

\subsection{一个超直观的小例子（不用扩散也能理解）}

令 \(\theta\in\mathbb{R}^2\)，输出 \(x_0\in\mathbb{R}^3\)：

\[
x_0 = A\theta,\quad r(x_0)=\tfrac12\|x_0\|^2
\]

其中 A 是 \(3\times2\)。

\begin{itemize}
  \item Jacobian：\(\frac{\partial x_0}{\partial\theta}=A（3×2）\)
  \item reward 对 \(x_0\) 的梯度：\(\nabla_{x_0} r = x_0\)（3维）
  \item 最终梯度：
  
  \[
  \nabla_\theta r = (\nabla_{x_0} r)^\top A = x_0^\top A
  \]
  
  这是一个 2 维向量。
\end{itemize}

注意：你只需要算 \(x_0^\top A\)，不需要把所有 \(J\) 展开成一个大对象再做运算（现实里 \(A\) 巨大时更是如此）。

\hrule

\subsection{VJP vs JVP（顺便区分一下）}

\begin{itemize}
  \item \textbf{VJP：} \(v^\top J\) ——反向模式（backprop）擅长，适合“标量输出、参数很多”（深度学习最常见）
  \item \textbf{JVP：} \(J u\) ——正向模式擅长，适合“输出很多、输入少”或做某些二阶/敏感度分析
\end{itemize}

我们训练网络一般都是标量 loss，所以几乎永远在做 \textbf{VJP}。

\hrule

\subsection{和你扩散训练的对应关系}

扩散里 \(x_0=g_\theta(\epsilon)\) 是一串迭代/网络调用的复合函数。你想要

\[
\nabla_\theta r(g_\theta(\epsilon))
\]

autograd 就是在沿着这串计算图做 VJP：把 \(\nabla_{x_0} r\) 作为“上游梯度”往回传，最后给你每个 UNet/LoRA 参数的梯度。

\hrule

如果你想，我也可以用你代码里的一段（比如 x\_0\_hat -> vae.decode -> online\_loss\_fn）具体指出：哪一行开始产生“上游梯度”，哪一行会因为 no\_grad()/detach() 把 VJP 链条切断。这样你就能清楚哪里能反传、哪里不能。
